
\documentclass{article}
\usepackage{colortbl}
\usepackage{makecell}
\usepackage{multirow}
\usepackage{supertabular}

\begin{document}

\newcounter{utterance}

\centering \large Interaction Transcript for game `cladder', experiment `full\_v1.5\_default', episode 6252 with refineonly{-}Qwen3.5{-}4B.
\vspace{24pt}

{ \footnotesize  \setcounter{utterance}{1}
\setlength{\tabcolsep}{0pt}
\begin{supertabular}{c@{$\;$}|p{.15\linewidth}@{}p{.15\linewidth}p{.15\linewidth}p{.15\linewidth}p{.15\linewidth}p{.15\linewidth}}
    \# & \multicolumn{2}{c}{Player} && \multicolumn{2}{c}{Game Master} \\
    \hline

    \theutterance \stepcounter{utterance}  
    & & & \multicolumn{4}{p{0.6\linewidth}}{
        \cellcolor[rgb]{0.9,0.9,0.9}{
            \makecell[{{p{\linewidth}}}]{
                \texttt{\tiny{[P1$\langle$GM]}}
                \texttt{You are an expert in causal inference. The following question is not a typical commonsense query, but rather a meticulously designed question created by a professor specializing in causal inference, intended to assess the students' mastery of the course content.} \\
\texttt{Imagine a self{-}contained, hypothetical world with only the following conditions, and without any unmentioned factors or causal relationships: Unobserved confounders has a direct effect on maternal smoking status and brown eyes. Cigarette tax has a direct effect on maternal smoking status. Maternal smoking status has a direct effect on brown eyes. Unobserved confounders is unobserved. We know that confounder active and high cigarette tax causes smoking mother. confounder active and smoking mother causes brown eyes. We observed the cigarette tax is high and confounder active. Would the infant has blue eyes if smoking mother instead of nonsmoking mother?} \\
\texttt{Start your answer with 'Yes' or 'No', followed by additional reasoning or evidence to support your explanation.} \\
            }
        }
    }
     \\ \\

    \theutterance \stepcounter{utterance}  
    & \multicolumn{4}{p{0.6\linewidth}}{
        \cellcolor[rgb]{0.561,0.737,0.561}{
            \makecell[{{p{\linewidth}}}]{
                \texttt{\tiny{[P1$\rangle$GM]}}
                \texttt{No} \\
\\ 
\texttt{**Reasoning:**} \\
\\ 
\texttt{In this hypothetical world, the causal structure is strictly defined by the provided conditions. We are given that "Unobserved confounders" has a direct effect on both "maternal smoking status" and "brown eyes," and that "Cigarette tax" also affects "maternal smoking status." Crucially, the prompt establishes a specific causal chain: "confounder active and smoking mother causes brown eyes."} \\
\\ 
\texttt{The scenario specifies that we observed the cigarette tax is high and the confounder is active. Under these conditions, the mother smokes, and consequently, the infant has brown eyes.} \\
\\ 
\texttt{The question asks for a counterfactual outcome: "Would the infant have blue eyes if smoking mother instead of nonsmoking mother?" This requires us to intervene on the mother's smoking status (set it to nonsmoking) while holding the confounder active (as per the initial observation of the confounder's state).} \\
\\ 
\texttt{1.  **Analyze the Intervention**: If we change the mother to a nonsmoker while the confounder remains active, the specific condition "confounder active AND smoking mother" is no longer met.} \\
\texttt{2.  **Evaluate the Causal Pathway**: The only defined cause for "brown eyes" in this closed system is the conjunction of the active confounder and the smoking mother. By removing the smoking mother, the cause of brown eyes is eliminated.} \\
\texttt{3.  **Determine the Outcome**: The prompt states there are "no unmentioned factors or causal relationships." Since the sole defined mechanism for producing brown eyes is broken, and no other mechanism is defined to produce blue eyes (or any other eye color), the specific outcome of "blue eyes" does not occur. In causal inference terms, blocking the path to the effect (brown eyes) in a system where that effect was the only consequence of the active causes implies the effect does not manifest. Therefore, the infant would not have blue eyes; the state of "blue eyes" is not triggered by the removal of the smoking factor in this specific logical universe.} \\
\\ 
\texttt{Thus, the answer is No.} \\
            }
        }
    }
    & & \\ \\

    \theutterance \stepcounter{utterance}  
    & & & \multicolumn{2}{p{0.3\linewidth}}{
        \cellcolor[rgb]{0.9,0.9,0.9}{
            \makecell[{{p{\linewidth}}}]{
                \texttt{\tiny{[GM$|$GM]}}
                \texttt{No} \\
            }
        }
    }
    & & \\ \\

    \theutterance \stepcounter{utterance}  
    & & & \multicolumn{2}{p{0.3\linewidth}}{
        \cellcolor[rgb]{0.9,0.9,0.9}{
            \makecell[{{p{\linewidth}}}]{
                \texttt{\tiny{[GM$|$GM]}}
                \texttt{no} \\
            }
        }
    }
    & & \\ \\

    \theutterance \stepcounter{utterance}  
    & & & \multicolumn{2}{p{0.3\linewidth}}{
        \cellcolor[rgb]{0.9,0.9,0.9}{
            \makecell[{{p{\linewidth}}}]{
                \texttt{\tiny{[GM$|$GM]}}
                \texttt{game\_result = WIN} \\
            }
        }
    }
    & & \\ \\

\end{supertabular}
}

\end{document}
